\subsection*{一步一步建立直观}

\paragraph{1. 把坐标解释成函数的系数}

在子空间 $\{a+bx:a,b\in\mathbb R\}$ 中，一个函数可以用系数对 $(a,b)$ 表示。右图的箭头现在代表系数，不代表物体在平面中移动。函数逐点相加对应系数相加，因此同一个二维图可以帮助理解这个函数空间。



\paragraph{2. 抽象来自保留运算规则}

把两份函数先相加再数乘，与分别数乘后相加，会在每一个输入点得到同一数值。这条分配律不依赖函数像不像箭头。向量空间的公理提取的正是这些可复用的运算关系，让几何语言能够进入函数、序列和多项式。



\paragraph{3. 选择基后才获得坐标图}

在上述二维子空间中，$1$ 与 $x$ 可作基；若只保留 $1+x$ 的倍数，就像把右图限制在一条线上。坐标图依赖所选有限维子空间，不能由此宣称所有函数都只有两份坐标。抽象定义比这个示意图的适用范围更广。
